% =====================================================================
%  seminar テーマ定義（既存 Marp デザインの Beamer 再現）
%  ・濃青の恒常ヘッダーバー（セクション名＋ページ番号）
%  ・専用表紙（アクセントライン＋グラデ背景）
%  ・本文は白地・濃青見出し
%  main.tex から \input する（パッケージではなく生プリアンブル）
% =====================================================================

% ---- ベース：装飾の少ないテーマから組む -----------------------------
\usetheme{default}
\usefonttheme{professionalfonts}
\setbeamertemplate{navigation symbols}{}   % ナビゲーション記号を消す

% ---- TikZ（ヘッダー・表紙・矢印などで使用）-------------------------
\usepackage{tikz}
\usetikzlibrary{arrows.meta, positioning, calc, shapes.geometric, fit, patterns}

% ---- 配色（ルール準拠）---------------------------------------------
\definecolor{HeaderBlue}{HTML}{0f5a7a}   % ヘッダー・見出しの濃青
\definecolor{DateGray}{HTML}{334b5e}     % 表紙の日付グレー
\definecolor{TitleBG}{HTML}{edf5fb}      % 表紙グラデの薄青

% コンテンツ用パレット（図・ハイライト・矢印で共通利用）
\definecolor{ArrowBlue}{HTML}{1F6FC4}    % 矢印・強調の青
\definecolor{HiPink}{HTML}{F4B6BE}       % ハイライト見出し（暖色）
\definecolor{HiBlue}{HTML}{C7DCEF}       % ハイライト見出し（寒色）
\definecolor{IpaBlue}{HTML}{1F5FC4}      % グラフ系列1（青）
\definecolor{DipeOrange}{HTML}{E8821E}   % グラフ系列2（橙）
\definecolor{SpRed}{HTML}{C0392B}        % 反応式：プロピレン強調の赤
\definecolor{SpBlue}{HTML}{1F6FC4}       % 反応式：エテン・エタン・CH0.5 の青

% 帯（ヘッダーバー）のパート別配色（寒色グラデ）。本文枠線は HeaderBlue のまま。
%  1-4 導入・全体像／5-9 反応・反応器／10-13 分離／14-17 最適化・経済・結言
\definecolor{barA}{HTML}{0F5A7A}   % 1-4  青緑
\definecolor{barB}{HTML}{1D6E6E}   % 5-9  シアン寄り
\definecolor{barC}{HTML}{2E4C8A}   % 10-13 青
\definecolor{barD}{HTML}{4A3C84}   % 14-17 紫

\setbeamercolor{normal text}{fg=black, bg=white}
\setbeamercolor{frametitle}{fg=HeaderBlue, bg=white}
\setbeamercolor{itemize item}{fg=HeaderBlue}
\setbeamercolor{itemize subitem}{fg=HeaderBlue}
\setbeamercolor{block title}{fg=white, bg=HeaderBlue}
\setbeamercolor{block body}{fg=black, bg=TitleBG}
\setbeamercolor{alerted text}{fg=HeaderBlue}

% ---- フォントサイズ（Marp px をおおよそ pt に対応）------------------
\setbeamerfont{frametitle}{size=\large, series=\bfseries}   % スライド見出し ~44px

% ---- 恒常ヘッダーバー（headline）-----------------------------------
%  左：セクション名（無ければロゴ枠）／右：ページ番号
\newcommand{\HeaderLogo}{}   % ロゴを入れるなら \renewcommand で画像指定
% 本線の総枚数（分母）。appendix に入ると分子がこれを超え、本線/付録を見分けられる。
\newcommand{\MainTotal}{17}
% パート別に帯色を選ぶ：1-4=barA／5-9=barB／10-13=barC／14-17=barD。
% スライド4の独自帯（CC・GCC）など headline を上書きする箇所でも呼べるよう独立コマンド化。
% 併せて本文アクセント色 HeaderBlue 自体をそのパート色へ再定義し、各スライドの
% 図・枠線・矢印・ボックス（HeaderBlue 参照）をパート色に連動させる（\colorlet は大域的）。
% headline はフレーム本文より前に組まれるため、ここでの再定義が本文に効く。
\newcommand{\SetBarColor}{%
  \ifnum\insertframenumber<5\setbeamercolor{headerbar}{fg=white, bg=barA}\colorlet{HeaderBlue}{barA}%
  \else\ifnum\insertframenumber<10\setbeamercolor{headerbar}{fg=white, bg=barB}\colorlet{HeaderBlue}{barB}%
  \else\ifnum\insertframenumber<14\setbeamercolor{headerbar}{fg=white, bg=barC}\colorlet{HeaderBlue}{barC}%
  \else\setbeamercolor{headerbar}{fg=white, bg=barD}\colorlet{HeaderBlue}{barD}%
  \fi\fi\fi}
\setbeamertemplate{headline}{%
  \ifnum\insertframenumber>0\relax
    \SetBarColor
    \begin{beamercolorbox}[wd=\paperwidth, ht=4.5ex, dp=1.8ex,
                           leftskip=1.2em, rightskip=1.2em]{headerbar}%
      \color{white}\bfseries\large
      \insertsectionhead%
      \hfill
      \insertframenumber/\MainTotal%
    \end{beamercolorbox}%
  \fi
}
\setbeamercolor{headerbar}{fg=white, bg=barA}
\setbeamertemplate{footline}{}             % フッターは使わない
\setbeamersize{text margin left=3mm, text margin right=3mm}  % 本文を横幅ぎりぎりまで

% 【必須ルール】帯と本文の間の余白を詰める（全スライド共通）。
%  beamer 既定の天井スキップで本文が帯から離れるのを、headline 直後の負スキップで打ち消す。
%  情報量最優先：無駄な上余白は作らない。各スライドは [t] 推奨。
\addtobeamertemplate{headline}{}{\vskip-2.6ex}

% frametitle はバーの直下に置く（左寄せ・濃青・太字）
\setbeamertemplate{frametitle}{%
  \vskip0.6ex
  \hspace{0.6em}\usebeamerfont{frametitle}\usebeamercolor[fg]{frametitle}\insertframetitle\par
  \vskip0.2ex
}

% ---- 箇条書きをシンプルに ------------------------------------------
\setbeamertemplate{itemize item}{\textbullet}
\setbeamertemplate{itemize subitem}{--}

% =====================================================================
%  再利用ヘルパー（どのスライドでも使える共通部品）
% =====================================================================
% 青い矢印 →（行中に置ける。結論や対応関係の明示に）
\newcommand{\bluearrow}{\tikz[baseline=-0.5ex]\draw[ArrowBlue, line width=1pt,
  -{Stealth[length=2.2mm]}] (0,0)--(0.55,0);}

% ハイライト見出し  \hilite{色}{テキスト}  例: \hilite{HiPink}{IPA選択率の温度依存性}
\newcommand{\hilite}[2]{\colorbox{#1}{\bfseries\normalsize #2}}

% パラメータ表  \paramtable{ 項目 & 値 \\ ... }  （左:項目 右:値、単位は数式推奨）
\newcommand{\paramtable}[1]{%
  {\footnotesize\renewcommand{\arraystretch}{1.15}%
   \begin{tabular}{@{\hspace{1em}}ll@{}}#1\end{tabular}}}

% =====================================================================
%  専用表紙コマンド  \seminartitlepage
%  ルール準拠：アクセントライン(120x8px相当) + 160度グラデ背景
%               h1 60px / h2 34px濃青 / 著者27px / 日付22pxグレー
%  使い方: \title{} \subtitle{} \author{} \date{} を設定後に呼ぶ
% =====================================================================
\newcommand{\seminartitlepage}{%
  \begin{frame}[plain]
    \begin{tikzpicture}[remember picture, overlay]
      % 背景グラデ（左上 薄青 → 右下 白、斜め）
      \shade[top color=TitleBG, bottom color=white, shading angle=20]
        (current page.south west) rectangle (current page.north east);
      % アクセントライン（上から少し下げて左に）
      \fill[HeaderBlue]
        ([xshift=9mm, yshift=-12mm]current page.north west)
        rectangle ++(16mm, 1.1mm);
      % テキストブロック（左寄せ・縦中央付近）
      \node[anchor=center, text width=0.86\paperwidth, align=center]
            at (current page.center) {%
        {\fontsize{22pt}{28pt}\selectfont\bfseries\color{black}\inserttitle\par}
        \vspace{0.3em}
        {\fontsize{16pt}{20pt}\selectfont\bfseries\color{HeaderBlue}\insertsubtitle\par}
        \vspace{1.0em}
        {\fontsize{13pt}{17pt}\selectfont\color{black}\insertauthor\par}
        \vspace{0.4em}
        {\fontsize{11pt}{14pt}\selectfont\color{DateGray}\insertdate\par}
      };
    \end{tikzpicture}
  \end{frame}
}
